%=============================================================================
\section*{日常速查表}
\addcontentsline{toc}{section}{日常速查表}
%=============================================================================

以下是你在集群上最常用的操作，不需要翻教程正文，直接复制粘贴即可。

\subsection*{连接与节点}

\begin{table}[h]
\centering
\small
\begin{tabular}{lp{9cm}}
\toprule
操作 & 命令 \\
\midrule
SSH 登录 & \texttt{ssh your\_utln@login-prod.pax.tufts.edu} \\
快速开一个 CPU 计算节点 & \texttt{srun -p batch -t 0-2:00:00 -n 2 --mem=4g --pty bash} \\
快速开一个 GPU 计算节点 & \texttt{srun -p gpu -t 0-2:00:00 -n 2 --mem=4g --gres=gpu:1 --pty bash} \\
退出计算节点 & \texttt{exit} \\
\bottomrule
\end{tabular}
\end{table}

\begin{tipbox}
很多操作（查目录大小、安装 Conda 环境等）在登录节点上会很慢，\textbf{开一个 CPU 计算节点再操作会快很多}。
\end{tipbox}

\subsection*{作业管理}

\begin{table}[h]
\centering
\small
\begin{tabular}{lp{9cm}}
\toprule
操作 & 命令 \\
\midrule
提交作业 & \texttt{sbatch script.sh} \\
查看我的作业 & \texttt{squeue --me} \\
取消某个作业 & \texttt{scancel JOBID} \\
取消我的所有作业 & \texttt{scancel -u \$USER} \\
查看已完成作业的效率 & \texttt{seff JOBID} \\
\bottomrule
\end{tabular}
\end{table}

\subsection*{资源与存储}

\begin{table}[h]
\centering
\small
\begin{tabular}{lp{9cm}}
\toprule
操作 & 命令 \\
\midrule
查看我的配额（实时刷新） & \texttt{watch -n 5 quota -s} \\
查看 Home 目录各文件夹大小 & \texttt{du -sh \textasciitilde/* \textasciitilde/.* 2>/dev/null | sort -rh | head -20} \\
查看集群 GPU 空闲情况 & \texttt{module load hpctools \&\& hpctools}（选 2） \\
查看各分区状态 & \texttt{sinfo} \\
查看 GPU 使用（在计算节点上） & \texttt{nvidia-smi} \\
\bottomrule
\end{tabular}
\end{table}

\subsection*{文件传输（在本地电脑上运行）}

\begin{table}[h]
\centering
\small
\begin{tabular}{lp{9cm}}
\toprule
操作 & 命令 \\
\midrule
上传项目目录 & \texttt{rsync -azP ./project/ user@login-prod.pax.tufts.edu:\textasciitilde/project/} \\
下载结果 & \texttt{rsync -azP user@login-prod.pax.tufts.edu:\textasciitilde/project/results/ ./results/} \\
\bottomrule
\end{tabular}
\end{table}

\subsection*{环境管理}

\begin{table}[h]
\centering
\small
\begin{tabular}{lp{9cm}}
\toprule
操作 & 命令 \\
\midrule
加载 Conda & \texttt{module load miniforge/25.3.0} \\
激活环境 & \texttt{conda activate rl\_env} \\
清理 uv 缓存（容易撑满配额） & \texttt{rm -rf \textasciitilde/.cache/uv} \\
清理 pip 缓存 & \texttt{rm -rf \textasciitilde/.cache/pip} \\
清理 conda 包缓存 & \texttt{conda clean --all -y} \\
\bottomrule
\end{tabular}
\end{table}

\newpage
